% Shared exercise chunk: l4-positionwise-fitting
\begin{exercisechunk}
\item \textbf{Fitting parity with independent-bit models.\suppmark}
\label[exercise]{ex:positionwise-fitting}
\exerciseprofile{2}{2}{2}{\exproof\quad\excalculation}{%
Compare fitting the parity law directly with fitting it using independent-bit
models.}
Let \(T\ge2\), let \(\cX=\{0,1\}^T\), and define
\[
 \operatorname{par}(x_{<T})
 =\left(\sum_{t=1}^{T-1}x_t\right)\bmod 2,
 \qquad x_{<T}\in\{0,1\}^{T-1}.
\]
For \(t\in[T]\), let \(X_t:\cX\to\{0,1\}\) be the \(t\)-th
coordinate map, and write \(X_{<T}=(X_1,\ldots,X_{T-1})\).
Let \(P_{\mathrm{par}}\) be the probability mass function on \(\cX\)
under which
the first \(T-1\) bits are independent fair coins and the last bit is their
parity. Let \(Q_{\mathrm{coin}}\) be the probability mass function on
\(\cX\)
under which all \(T\) bits are independent fair coins. Fix \(n\ge1\) and
points \(x^{(1)},\ldots,x^{(n)}\in\cX\) satisfying
\(P_{\mathrm{par}}(x^{(i)})>0\) for every \(i\in[n]\). For any probability
mass function \(Q\) on \(\cX\), define its empirical log loss per token on
these points by
\[
\widehat L_n(Q)=-\frac1{nT}\sum_{i=1}^n\log Q(x^{(i)}).
\]
\begin{enumerate}[(a),beginpenalty=10000]
\item Show that
\[
 \widehat L_n(P_{\mathrm{par}})=\frac{T-1}{T}\log 2,
 \qquad
 \widehat L_n(Q_{\mathrm{coin}})=\log 2,
 \qquad
 \widehat L_n(Q_{\mathrm{coin}})-\widehat L_n(P_{\mathrm{par}})
 =\frac{\log 2}{T}.
\]
Verify that for every \(x_{<T}\in\{0,1\}^{T-1}\),
\[
 P_{\mathrm{par}}\!\left(
 X_T=\operatorname{par}(x_{<T})\mid X_{<T}=x_{<T}\right)=1,
 \qquad
 Q_{\mathrm{coin}}\!\left(
 X_T=\operatorname{par}(x_{<T})\mid X_{<T}=x_{<T}\right)=\frac12.
\]
\item Let \(\cQ_{\mathrm{ind}}\) be the class of independent-bit laws
\[
 Q_q=\bigotimes_{t=1}^T\operatorname{Bernoulli}(q_t),
 \qquad q=(q_1,\ldots,q_T)\in[0,1]^T.
\]
First verify that \(Q_{\mathrm{coin}}=Q_{(1/2,\ldots,1/2)}\) belongs to
\(\cQ_{\mathrm{ind}}\). Then derive
\[
 \E_{P_{\mathrm{par}}}[-\log Q_q]
 =\sum_{t=1}^T\left[-\frac12\log q_t-\frac12\log(1-q_t)\right]
\]
and show that the unique minimizer of the population sequence log loss over
\(\cQ_{\mathrm{ind}}\) is \(Q_{\mathrm{coin}}\).
\end{enumerate}

\paragraph{Interpretation.}
Part~(a) shows that one completely missed final-position dependency contributes
only \((\log 2)/T\) to average per-token log loss, although the independent
model assigns only probability \(1/2\) to the correct final bit. Part~(b)
identifies the role of model
misspecification: the best independent-bit law matches every one-coordinate
marginal and is \(Q_{\mathrm{coin}}\), but no law in this class represents the
parity dependency. Thus optimal fitting within a factorized class need not
recover a sequence-level dependency.
\end{exercisechunk}
